\chapter{Netz-, Material- und Randverträge}
\label{api:mesh}\index{Mesh}\index{Material}\index{Dirichlet}
\begingroup\raggedright
\section{Mesh und Topologie}
\texttt{Mesh} besitzt endliche reelle \texttt{points[n,3]} in m;
\texttt{tetra4[e,4]}, \texttt{tetra10[e,10]}, \texttt{triangle3[f,3]} und
\texttt{triangle6[f,6]} sind Integer-Konnektivitäten. Optionale IDs/Tags haben
je Entitätszahl eine Zeile, Sätze enthalten eindeutige nullbasierte Indizes.
Eingaben werden kopiert; Arrays/Satzzuordnungen sind unveränderlich. Leere
optionale Zellblöcke sind zulässig, doppelte/topologisch ungültige Entitäten nicht.
Solver verlangen passende aktive Volumenzellen; gemischte Elementordnungen
sind ausgeschlossen. Geometrische Entartung wird in den Elementen/Importern
geprüft. \texttt{inspect\_boundary\_topology} bindet Außenflächen an Zellen,
bewahrt Satzmitgliedschaft und diagnostiziert innenliegende, nichtmannigfaltige
oder falsch gebundene Flächen. BoundaryFace/BoundaryTopology sind Daten des
Topologieberichts; Einheiten und Arraytypen stehen im Inventar.

Geradkantige Tet10 verwenden feste Kantenreihenfolge und mittige Zwischenknoten;
keine gekrümmte Geometrie. \texttt{Mesh.cells}, \texttt{cell\_ids},
\texttt{n\_cells}, \texttt{element\_order} vereinheitlichen den aktiven Zelltyp.
Tiefe Geometrie-/Quadratur-/Assembly-Bausteine sind Forschungs-API, auch wenn
Anwendungsbeispiele sie zur Lehre importieren. Mindestqualität, Orientierung
und Außen-Normalen gelten gemäß Theorieband und Doc/40--42.
Minimalbeispiel: \nolinkurl{examples/steady_heat.py}; Import und Datenformen
werden durch \nolinkurl{tests/unit/test_mesh_io.py} geprüft.

\section{Materialtensoren und Regionen}
\texttt{ThermalMaterial}: symmetrische positiv definite Leitfähigkeit $k[3,3]$
in W/(m K), optional positive Dichte kg/m$^3$ und Wärmekapazität J/(kg K) für
transiente Modelle. Quellen-/Katalogmetadaten beschreiben konstante Werte und
deren Gültigkeit; sie erzeugen kein temperaturabhängiges Gesetz.
\texttt{ElasticMaterial}: $C[6,6]$ in Pa, symmetrisch positiv definit;
\texttt{isotropic(young, poisson)} prüft $E>0$, $-1<\nu<1/2$.
\texttt{DielectricMaterial}: $\kappa[3,3]$ in F/m, SPD.
\texttt{PiezoelectricMaterial}: Stress-Charge-Form $C[6,6]$, $e[3,6]$
in C/m$^2$, $\kappa[3,3]$; Rotation/Data-Sheet-Konversion behandelt
Voigt-/Schub-/Permittivitätsbedingungen ausdrücklich. Eine strain-charge-
Datenblattpermittivität bei konstanter Spannung wird nicht ungeprüft als
stress-charge-Permittivität übernommen. Dichte ist für Dynamik erforderlich.
\texttt{ThermalExpansion}: $\alpha[6]$ in 1/K und absolute Referenztemperatur
in K, isotrope/principale Konstruktoren und Achsrotation gemäß Doc/18.

Homogenes Material oder genau ein Objekt pro Eingabezelle. Regionen müssen
vollständig/eindeutig zugeordnet sein; sich widersprechende Überlappungen und
Lücken erzeugen \texttt{MaterialAssignmentError}. \texttt{assign\_region\_materials}
liefert MaterialAssignment in Zellreihenfolge mit Zuordnungsprovenienz.
\texttt{Model} konvertiert nur erlaubte Einheiten; Direktobjekte erwarten SI.
Arrays sind owned/readonly; Quelle und Achsen bleiben prüfbare Metadaten.
Minimalbeispiele: \nolinkurl{examples/piezo_d33.py},
\nolinkurl{examples/thermal_expansion.py}; vollständige tensorielle
Konstruktoren stehen in den Signaturen/Docstrings des Inventars.

\section{Dirichlet, Elektroden und Lasten}
\texttt{Dirichlet(dofs, values)} bindet globale Feld-DOF-Indizes an SI-Werte.
Konflikte werden abgelehnt. Für $u$: m, $T$: K, $\phi$: V.
\texttt{Electrode} wählt Außen-Dreiecksflächen: festes Potential in V oder
schwebende Gesamtladung in C. Ausgedehnte Elektrode ist equipotential;
überlappende Konflikte und ungebundene Innenflächen sind Fehler. Schwebende
Elektroden benötigen weiterhin eine elektrische Referenz; Elektrode ist
kein beliebiger Ein-Knoten-Potentialzwang. ElectrodeState beschreibt
intern die Kondensation, ohne eigene stabile Anwenderzusage.

Lasten: volumetrische Wärmequelle W/m$^3$, äußerer Wärmestrom W/m$^2$,
Konvektion W/(m$^2$ K) mit Umgebung K; Volumenkraft N/m$^3$, Traktion N/m$^2$,
Punktkraft N; Volumenladung C/m$^3$, Oberflächenladung C/m$^2$.
Mappings adressieren Außenflächen-Zeilen oder globale Knoten gemäß Signatur.
Callbacks erhalten physische Punkte in m und liefern endliche SI-Werte;
Assembly friert effektive diskrete Lasten in Provenienz ein, keinen Python-
Quellcode eines Callbacks. \texttt{Convection} ist ein
Lastdatensatz des Wärmevertrags; niedrige Lastintegrationsfunktionen
bleiben Forschungsbausteine. Beispiel: \nolinkurl{examples/thermopiezoelectric/run.py}.
\endgroup
